\section{Conclusion}
\label{sec:conclusion}

We introduced \method{}, a system that turns a long-form LLM answer into an
interactive visual image: a generated picture overlaid with clickable hotspots,
where each region opens its own explanation and follow-up thread. The central
lesson is practical: if the image generator is not guaranteed to obey a layout,
the system should not trust the layout as the final interaction layer. It
should generate first, then ground the visible result. We implemented this idea
in an open-source, provider-agnostic system spanning infographic, map, and
scene modes, and evaluated it with completion metrics, source attribution for
hotspot alignment, and a strict visual-alignment gate.

\paragraph{First work in this direction.}
Converting an LLM answer into a navigable, region-interactive visual artifact
with vision-grounded hotspots is, as far as we know, a new task not addressed
directly by prior work in text-to-image generation, visual grounding,
interactive visualization, or structured LLM output. Our contribution is not a
new foundation model, but a system pattern: structure the answer, render it,
ground the rendered regions, and attach local interaction to each region.

\paragraph{Application prospects.}
The interactive visual answer format applies naturally to domains where long
answers already contain visual or spatial structure. In
\textbf{education}, a student can explore a complex process (e.g., the HTTP
rendering pipeline) by clicking the stage they care about rather than reading
an entire explanation. In \textbf{technical documentation and knowledge
management}, answers become navigable artifacts where each region supports
targeted follow-up, reducing re-reading. In \textbf{tourism and spatial
guidance}, hand-drawn maps with explorable landmarks and routes offer a more
intuitive interface than text directions. In \textbf{data communication},
infographic answers make comparisons and funnels immediately scannable while
preserving drill-down detail. As text-to-image and visual-grounding models
continue to improve, the fidelity and alignment of such interactive visual
answers should improve as well, widening the range of questions that benefit
from this format.

\paragraph{Limitations and future work.}
Current limitations are mostly where the pipeline meets imperfect vision
models. Dense infographics, small objects, and complex maps can still produce
weak or partial groundings. Click targets also remain rectangular in the normal
pipeline, even when masks provide a better-looking preview; polygon-level hit
testing would better fit irregular regions. File context is still text-oriented,
so richer document parsing remains future work. Promising directions include
shareable HTML exports, stronger OCR-aware grounding, richer mask-based
interaction, automated visual QA for hotspot accuracy, and synchronized
cross-device history. The open-source release and benchmark are intended to
make this format easier to inspect, reproduce, and extend.

\paragraph{Open source.}
Code, demos, a live project page, and a technical report are publicly
available.\footnote{\url{https://github.com/wencanjiang/ChatImage}}
